\section{Related Work}

Automatic passenger counting is a foundational capability in public transit. While video-based \ac{APC} systems deploy cameras at vehicle doors for neural tracking \cite{pronello2023evaluating}, they require dedicated hardware and raise privacy concerns. In contrast, Wi-Fi sensing passively monitors radio emissions from passenger mobile devices at zero incremental hardware expense, offering an attractive, privacy-preserving alternative.

Wi-Fi sensing operates primarily on two signal representations. Deep learning has accelerated Wi-Fi human sensing \cite{ahmad2024wifisensing}. Channel State Information (CSI) provides subcarrier-level resolution for passenger counting in vehicle cabins \cite{wang2024csi}, but requires specialized firmware modifications unavailable on standard routers. The \ac{RSSI}, by contrast, is ubiquitous across off-the-shelf hardware. RSSI fingerprinting is widely used for localization \cite{singh2021rssisurvey}, probe-based occupancy counting \cite{nitti2020iabacus}, signature tracking \cite{myrvoll2017counting}, and ambient activity classification \cite{sigg2014telepathic}.

Integrated sensing and communications leverages telecommunications infrastructure as dual-functional sensing nodes \cite{liu2022isac,gonzalezprelcic2024isac}. Machine learning plays a central role in the ISAC processing pipeline \cite{luong2025mlisac}, extracting environment states from ambient signals. Our work adopts this opportunistic paradigm, classifying passenger transit states from \ac{RSSI} time series using standard onboard access points without dedicated radar units.

Sequence modeling has advanced through selective state spaces such as Mamba \cite{gu2023mamba} and Mamba-3 \cite{gu2026mamba3}. In parallel, self-supervised joint-embedding architectures (I-JEPA \cite{assran2023ijepa}, LeJEPA \cite{balestriero2025lejepa}) learn robust representations by predicting in latent space. Recent extensions adapt JEPAs to tabular and temporal data \cite{thimonier2024tjepa,tsjepa2024,cfjepa2026}, offering promise for Wi-Fi sensing where labeled traces are costly \cite{radwan2025sslwifi}. While prior work on this task relied on classical Gaussian Processes (0.756 \ac{MCC} \cite{ribeiro2026eucnc}), this paper is the first to deploy self-supervised world models for opportunistic Wi-Fi sensing in public transit.
